\section{Energy}
\label{appendix:energy}

In this Appendix, we present the modelling of the energy sector, globally and in the UK. 

Our approach is to take prices mostly from \citet{wayEmpiricallyGroundedTechnology2022}, and quantities from various scenarios (UK Climate Change Committee scenarios, \citet{wayEmpiricallyGroundedTechnology2022}, and our own). In Section \ref{app_energy_prices}, we describe how we obtain prices. In Section \ref{app_energy_mix}, we describe how we construct the quantities under various scenarios. We present the results in Section \ref{app_energy_results}.

Before diving into this, it is useful to introduce some key concepts. First, there is a distinction between final and useful energy: final energy is the processed form delivered to consumers and firms, while useful energy is what remains after accounting for conversion losses, representing the actual service (heat, kinetic energy, light, etc.). Second, the electricity sector must ensure market clearing at all times, ensuring supply meets demand. Additional solutions must be implemented in high renewables systems due to the inherent intermittence—variability in supply—of solar and wind. Storage captures excess energy for later use, and flexibility balances supply and demand through mechanisms like demand response and grid interconnections. Finally, the price of energy in the market is set by the most expensive production unit needed to meet demand at a given time, known as the marginal price. It results from a bidding process and reflects both capital expenditure (capex) for infrastructure and operational expenditure (opex) for ongoing costs like fuel and maintenance. In our exercise,we assume that, on average, the electricity market clearing price is close to the average price \FL{cost?} of electricity.




\subsection{Prices}
\label{app_energy_prices}

All prices are expressed in USD 2020/MWh. 

\paragraph{Main sources} When available, we take the price time series from \cite{wayEmpiricallyGroundedTechnology2022} (solar PV, wind, oil, gas, solid (coal), hydrogen (P2X), batteries) and cost assumptions from BEIS and DENSZ in the UK for the other energies. \cite{wayEmpiricallyGroundedTechnology2022} forecast prices using Wright's and Moore's laws, fitting it to past data. It shows a great decrease in cost -- as experience accumulates -- for solar, wind, hydrogen and batteries but no trend for nuclear and fossil fuel based technologies. We are not using \cite{wayEmpiricallyGroundedTechnology2022} gas generated electricity (50 USD 2020/MWh) because it doesn't allow for proper comparison with Gas+CCS electricity generation, and we cannot decompose it to get the cost of hydrogen electricity generation (by adding H2 cost -- with a conversion factor of 0.5 -- to the gas plant fixed cost (14 USD/MWh). See Table \ref{tab_sources_prices_energy} for details.

\begin{figure}
    \centering
    \caption{Evolution of the cost of energy (constant USD 2020/MWh)}
    \label{g_price_energy}
     \includegraphics[width=\textwidth]{figures/figures_energy/g_price_energy.jpeg}
     {\\\raggedright\footnotesize\textbf{Note}: Refer to  Table \ref{tab_sources_prices_energy} for details. Panels are ordered in descending order relative to the 2050 cost of a MWh, from biomass+CCS electricity at 280.9 USD/MWh to Solar PV electricity 15.6 USD/MWh (including network cost). \\}
\end{figure}

We consider the price of technologies in BEIS as ``levelised cost for a plant built in 2025'' as it show little trend of decreasing over time. Indeed carbon capture and storage technologies existing cost data show no decreasing trends: see \cite{wayEmpiricallyGroundedTechnology2022}, appendix 1.6.1 and \cite{bacilieriAssessingRelativeCosts2023}, finding no evidence of declining cost of CCS.

\paragraph{Supplementary assumptions} Electricity generation cost include self-use of energy, included in the 0.5 conversion factor. However, we do not consider transmission losses in the network, nor losses in storage (P2X or batteries), 

\paragraph{Network cost} The transition of the energy systems requires investment in the electricity distribution network. First, because massive electrification means a scale up of the quantity of electricity that needs to be transmitted and thus the number of power lines, transformers, etc. Second, because the use of renewable electricity generation means that the decentralised production of electricity must be matched with the demand throughout the country. The grid must be able to accommodate such flexibility. \cite{wayEmpiricallyGroundedTechnology2022} model the two component independently for the global energy system. We take a simpler approach and consider a network cost of 5£ per MWh \citep{carbonbriefIndepthWholeSystem2017}. We assume it applies linearly to any additional unit of renewable electricity (Wind and solar) above 2022 electricity production of wind and solar. They account for about -550bp over 2022-2050 (in the Balanced Net Zero Scenario).\footnote{
%
A back of the envelope calculation gives a similar order of magnitude. Wind is about 20 USD/MWh in 2050. A 7 USD 2020/MWh is an increase of about 25\%, if applied to all renewables, 40\% of final energy in 2050 in Balanced Net Zero Scenario, the impact on the average price is about +10\%. With $\zeta=0.12$, the productivity impact would be 800bp.
%
} 
 
\begin{landscape}
\input{tables/tab_sources_prices_energy.tex}
\end{landscape}

\paragraph{Implied CCS cost} We can deduce from the comparison of gas and gas+ccs electricity generation cost the implied cost of CCS. The cost per MWh -- exclusive of carbon costs -- are respectively 53 and 81 £2018/MWh for gas without and with CCS. (table 4.3 \cite{beisElectricityGenerationCost2020},"CCGT H
Class"	vs "CCGT + CCS Post Combustion (FOAK)"). Assuming an average 469gCO2/kWh per gas-generated electricity (IPCC, \cite{edenhoferSpecialReportRenewable2011}, Appendix II Methodology, A.II.4, p982), it means that the cost surplus of CCS (28 £2018/MWh) is equivalent to a cost of CCS of 59.7 £2018/tCO2, that is 81.8 USD 2020/tCO2. 


\paragraph{Macroeconomic indicators} We use UK GDP deflator to convert values into constant 2020£ (Table \ref{tab_inflation}), then convert them into constant USD 2020 (Table \ref{tab_echange_rate})

\begin{table}[H]
%\centering
\caption{UK GDP deflators}
\label{tab_inflation}
\begin{tabular}{@{}lll@{}}
\toprule
Report    & Unit used & GDP deflator to £2020\\ \midrule
BEIS 2016 & 2014£     & 1.133117                                             \\
BEIS 2020 & 2018£     & 1.067278                                             \\
BEIS 2023 & 2021£     & 0.996194                                             \\ \bottomrule
\end{tabular}
  {\\\raggedright \textit{\textbf{Source}: ONS, 2024, \href{https://www.ons.gov.uk/economy/grossdomesticproductgdp/datasets/uksecondestimateofgdpdatatables}{ GDP data tables}. Table A1 "Worksheet A1: National accounts aggregates". \\}}
\end{table}

\begin{table}[H]
%\centering
\caption{Dollar-Sterling exchange rate}
\label{tab_echange_rate}
\begin{tabular}{@{}ll@{}}
\toprule
Annual average   Spot exchange rate, US \$ into Sterling & XUAAUSS \\ \midrule
2020                                                     & 1.2837  \\ \bottomrule
\end{tabular}
 {\\\centering\raggedright \textit{\textbf{Source}: Bank of England, \href{https://www.bankofengland.co.uk/boeapps/database/fromshowcolumns.asp?Travel=NIxAZxSUx&FromSeries=1&ToSeries=50&DAT=RNG&FD=1&FM=Jan&FY=2019&TD=28&TM=Feb&TY=2024&FNY=Y&CSVF=TT&html.x=66&html.y=26&SeriesCodes=XUAAUSS&UsingCodes=Y&Filter=N&title=XUAAUSS&VPD=Y}{XUAAUSS time series}, 2024, \\}}
\end{table}


\citet[table 2.1, p.51]{cccDeliveringReliableDecarbonised2023} gives cost assumptions for batteries and other storage solutions; for instance, batteries in 2035 are assumed to cost 260-980 £2012/kW. All other solutions are even more expensive: pumps 1290-1830, compressed air 860-1420, liquid air 900-1050£/kW. We are not using these numbers as they require non trivial assumptions to convert into kWh.

%\paragraph{Alternative solution for storage and network cost} [\textit{not implemented}] One solution could be to consider as in WIMF that there is no electricity from hydrogen, and just add a mark up of 10-20£/MWh of renewable which includes hydrogen based electricity. And then apply a learning curve to this top up cost: with a learning curve between 0 (as in transmission losses or grid development) and the learning rate of batteries (which is the quickest learning technology in this bucket).





\subsection{Energy mix}
\label{app_energy_mix}

We first detail the underlying energy mix of the CCC and how to match it to WIMF, we then focus on hydrogen and storage. Finally, we discuss how to hybridise the scenarios.


\subsubsection{CCC energy mix}

We consider the energy mix of the CCC 6th Carbon budget \citep{cccSixthCarbonBudget2020, cccMethodologyReportSixth2020}.
The final energy demand is from the sheet "Scenario key metrics", the electricity generation mix from the sheet "Electricity generation data". 

\paragraph{Excluded rows.}
We exclude for the calculations the "Interconnection" as i) we do not know which technology is used, hence its price, ii) it does not correspond to actual production when negative, which is the case for Balanced Net Zero between 2023 and 2050.\footnote{\cite{cccMethodologyReportSixth2020}, p156, ``Interconnectors''.} We also do not consider ``Used curtailment'' for hydrogen production: it is surplus electricity production that needs to be used to ensure the stability of the system. Its marginal cost is then null, and if not used for hydrogen selection (or battery filling) it would be disposed of at a non-zero cost. Without further assumptions, its macro impact is unclear. 

\paragraph{Matching of energy sources.} We consider that ``Solid fuel'' corresponds to coal. ``Electricity generation: Thermal low-carbon'' corresponds to biomass and geothermal, for the sake of simplicity we assume it is biomass generated electricity. ``Electricity generation : Other renewable'' includes hydro, tidal and wave, we assume it is all hydro-electricity (limited in the UK).  ``Electricity generation : Other'' includes includes coal, gas CHP, reciprocating engines. For the sake of simplicity we consider it gas (CHP plant whereas the other gas power plants are CCGT).\footnote{Sensitivity shows that accounting for it as coal does not change much.}

\paragraph{Primary \textit{vs} final energy} Note that the primary energy consumption of the CCC includes the quantities used for electricity production. We want to use the final energy mix, that is that the gas used to produce electricity is not accounted for, but the generated electricity is. From the gas demand, we subtract the quantity of electricity produced by gas -- unabated gas, CCS gas and CHP gas ---, we subtract biomass electricity (with and without CCS) to bioenergy, and hydrogen generated electricity are subtracted to hydrogen. We assume a conversion factor of 0.5 for gas and hydrogen, and 0.4 for coal following \cite{wayEmpiricallyGroundedTechnology2022}\footnote{SI 2.6, p33.}, meaning that we divide the quantity of electricity by 0.5 (or 0.4) to get the required primary energy quantity. We assume the same yield for biomass than for coal. 

UK primary energy equivalent consumption in 2022 is 170Mtoe (ECUK, BEIS, 2022)\footnote{\href{https://www.gov.uk/government/statistics/energy-consumption-in-the-uk-2022}{Energy consumption in the UK 2022, ECUK 2022: Consumption data tables}}. The sum of the corrected quantities of energy in the Balanced Net Zero Pathway is 1866 TWh in 2022, which is about 160Mtoe (errors in conversion factors, density of energy, convention about curtailment and storage easily explain the 6\% gap).

\subsubsection{Electricity from hydrogen}

P2X fuel, standing for Power-to-X, plays a critical role in the net zero transition of energy systems by converting electricity from wind and solar to chemicals: hydrogen, ammonia, methane, methanol. These are used as long-term storage and transport means of energy in the form of hydrogen-based fuels or can be used directly for hard-to-electrify industrial activities. In the transition, hydrogen is the main P2X fuel, and we will consider it exclusively in this analysis.

 The storage target of hydrogen for electricity generation in 2070 is a share $\theta$ of final electricity generated annually from wind and solar: either a week (slow transition), a month (fast transition), or nothing (no transition) \citep{wayEmpiricallyGroundedTechnology2022}.\footnote{\label{footnote_energy_produced}Note that in WIMF, $\theta$ is expressed in final energy delivered from wind and solar, and not in the total energy generated by these technologies. The former being about 50\% higher than the latter. Indeed, renewables are producing even when not needed, this surplus energy is then either stored in batteries, converted to H2, traded, or disposed of. Is it the same for batteries capacities computation.}

\cite{wayEmpiricallyGroundedTechnology2022} consider that there is no electricity from hydrogen before 2042. We approximate the growth in P2X capacities for electricity generation as a linear trend between 0 in 2042 and the target $\theta=8.3\%$ in 2070 (in the high renewable "Fast transition" scenario) reaching 2.6\% in 2050.

\paragraph{WIMF scenario} To ease comparison, we explicitely consider storage in WIMF scenario. We add electricity from hydrogen in the Fast transition scenario: from 0 in 2042 to 8.3\% (1/12) in 2070 with a linear growth between 2040 and 2070, reaching 2.6\% of electricity production in 2050. It represents  8.3Ej of electricity from P2X fuel, with a 0.5 conversion factor it amounts to 16.7EJ of hydrogen (way below P2X production, ensuring it makes sense). P2X electricity is then on top of that and does not appear as such in figure 6 of the main text.\footnote{To check this computation, one could compare it to Figure S8 (I, L) \citep{wayEmpiricallyGroundedTechnology2022} and S9, p57 of the supplementary material, show that P2X production is limited before 2040, and P2X for electricity begins in 2042, to reach about 50EJ/year in 2070, that is about 8\% of the 630EJ of final energy consumption in 2070. Final use of hydrogen for the sectors (not used for electricity) is about 80EJ in 2050 and 140EJ in 2070. (corresponds to the P2X fuel in the fast transition scenario: 83 EJ in 2050 and 143 in 2070).}

\paragraph{CCC scenario} Electricity from hydrogen is already included in the CCC scenarios. In the 5 CCC scenario, electricity from hydrogen represents between 0.1\% (Headwinds
) and 3.0\% (Widespread Engagement) in 2050. The share of electricity being produced with hydrogen is in line with the \cite{wayEmpiricallyGroundedTechnology2022} assumptions that in a low renewable world, the reserve of P2X fuels, must cover 1 week worth of electricity, that is about 1.9\% of annual consumption on average. In a very high renewable scenario, as in the Fast Transition, there must enough to cover one month worth of electricity, that is about 8.3\% in 2070.\footnote{section SI3.3.3 of \cite{wayEmpiricallyGroundedTechnology2022} supplementary material}   However, the quantity of hydrogen produced in 2050 in the CCC scenarios is huge, between 161 and 376 TWh, and is probably overestimated given the residual consumption of gas in 2050. 


\subsubsection{Electricity Storage}

Electrifying the usages and increasing the share of renewable and intermittent energy (wind and solar) means increased needs for flexibility. Long-term storage is provided by P2X, whereas short-term daily storage is provided by Li-ion batteries and medium term multi-day storage is provided by flow batteries, namely Vanadium redox flow (VRF) (see appendix in \cite{wayEmpiricallyGroundedTechnology2022}, sections 3.1 and 3.2, for a review of the literature and the assumptions).

The need for battery capacity in a high renewable world, is estimated to be 20\% of total energy produced by wind and solar for daily usage (about the equivalent of 6h of final electricity demand), and multi-day capacity is estimated at 10\% of total energy produced by wind and solar. We follow \cite{wayEmpiricallyGroundedTechnology2022} by distinguishing between final electricity delivered and energy generated by wind and solar, the latter being as 50\% greater than the former, the excess producing hydrogen through electrolysis or being stored in batteries. 

In WIMF, short term storage using Li-ion batteries have a steap learning curve, with a Moore's law coefficient of $\mu=0.146$ (Table S26, p129) such as $c_t=e^{\ln(c_{t-1})\mu}$, a starting cost at 310\$/kWh in 2020, then decreasing to 4.5\$/kWh in 2050. While medium term VRF batteries have starting cost at 310\$/kWh and a Moore exponent $\mu=0.076$ ,a starting cost at 400\$/kWh in 2020, then decreasing to 44.4\$/kWh in 2050. The cost plotted in WIMF, figure 6.C correspond to the capacity-weighted average of short and medium term batteries, to obtain a cost of about 17\$/kWh in 2050.

Storage needs are expressed in capacities, to convert it into annual contribution to the energy "production", we multiply the daily capacity by 365, assuming that all the energy stored into both short term and medium term batteries are fully used everyday. The implication being that our results err on the higher assumptions for storage overall cost when comparing scenarios. We then consider annual "production" of energy via storage, dividing the capacity price by 365 to keep the overall cost consistent. 

\paragraph{WIMF methods} To compute the need for storage, WIMF distinguish between short term and medium term. The capacity for storage is $q_X$ with $X\in (ST,MT)$, the storage requirement is expressed as a share $\beta_X$ of generated energy from renewables ($q_R$). As described earlier, it is 20\% for short term storage, and 10\% for medium term storage in a high renewable world (0.0001 for both $\beta_{MT}$ and $\beta_{ST}$ in a low renewable world).  
\[q_{X}=\beta_X \frac{q_R}{365}.\] 
We add up the storage requirements of short term and medium term, we are considering them together, and express the storage capacity in batteries as $q_B= \beta_B\frac{q_R}{365}$ with $\beta_B=\beta_{ST}+\beta_{MT}=0.3$. In any given scenario, we can compute $q_R$ as 150\% of the energy delivered by wind and solar (by assumption). Note that the capacity is expressed in TWh but is daily or multi-day storage.



\paragraph{CCC scenario} Most of the CCC scenario have wind and solar generate between 71 and 85\% of the final electricity, that is between 25\% and 49\% of final energy. It is not a "high renewable, high electrification" scenario and thus does not require a lot of storage. Storage amounts to 0.06\% of total final energy in 2020 and 0.35\% in 2050, with a peak at 0.49\% in 2040. It represents up to 1.4\% of total electricity produced. Much less than the WIMF scenarios. Storage is already expressed in annual contribution. The underlying assumption is that daily use is constant. 

We compute the required capacities of batteries ($q_B$) following WIMF method in a high renewable world. We multiply the result by 365 to get an annual estimate (assuming that the whole capacity of batteries is used everyday) to compare to the CCC forecast. WIMF method gives a daily storage need of 0.098TWh in 2020 and 0.063TWh in 2050. To be compared to Balanced Net Zero: in 2022,  0.003TWh, and in 2050, 0.012TWh.

The impact of battery storage on productivity should then be very limited: daily 0.012 TWh in 2050, it is 4.34 annual TWh, it represents 0.35\% of final energy use (1245TWh). 

With an average price of energy of 37.2USD 2020/MWh in 2050, and a 2050 price of battery of 0.04USD 2020/kWh (43.8\$/MWh), the impact on price is 0.15 USD/MWh.\footnote{The price of a TWh of battery capacity in 2050 is 16,000USD/MWh, to be divided by 365 to get the price per day.} There is a very limited impact of batteries on productivity growth, a few dozen bp at max, 1bp on cumulative productivity in the Balanced Net Zero scenario. 


\subsection{Building mix scenarios}
\label{section_UK_WIMF_scenario}

We cannot apply directly the WIMF "energy mix" scenario to the UK, as the initial energy mix are not compatible (see Table \ref{tab_initial_price_energy}). We then adapt the WIMF scenario to the UK in three steps. 

\paragraph{Total quantity of energy} To downscale WIMF Fast Transition to the UK, we use the CCC assumptions about final energy consumption. However, WIMF scenarios assume a massive electrification of the economy. Thus, we cannot directly use the final energy quantity from CCC scenarios are electrification get rid of conversion losses. We consider of about 0.5 between fossil energy and electricity.\footnote{0.4 for biomass and coal, WIMF assumption.}. 
To compare the same quantity of final energy, we convert the WIMF energy into "final energy in electricity equivalent", multiplying fossil fuels by the 0.5 conversion factor. Note that this is not exactly the useful energy metrics defined by WIMF, as the useful energy also integrates final use conversion losses (and not only electricity production conversion losses). 

Our assumption is then that the CCC scenarios give an accurate forecast of "final energy in electricity equivalent" in 2050 expressed in TWh: TWh of oil, solid, gas, hydrogen, waste and bioenergy are multiplied by 0.5. 

\paragraph{Energy mix} We assume that the UK will reproduce the WIMF Fast transition energy mix in 2050. We multiply the total "final energy in electricity equivalent" of the CCC scenario by the "Fast transition" WIMF energy shares. We convert it back to final energy (oil, solid, gas, hydrogen, waste and bioenergy quantities are divided by 0.5). 

CCC scenario show little change in the energy mix from 2022 to 2025, and most projects are already underway. Most of the new technologies, such as gas+CCS for electricity generation are included in the CCC energy mix starting 2026 and 2027. We assume that CCC got the trajectory right until 2025. 

We create a trajectory to match the WIMF Fast transition global energy mix in 2050 starting from the 2025 Balanced Net Zero UK energy mix.
Between the initial quantity of energy for each source in 2025 and the new production target to match WIMF scenario, we assume a linear growth. 

\paragraph{Solar and Wind correction} 
The Fast transition uses a large quantity of solar energy, more than what is available in the UK, where wind is an easier solution. We correct this factor, by making sure that the 2050 ratio between wind and solar reflects the Balanced Net Zero pathway the CCC (ratio of 5 between offshore + onshore wind and solar). We attribute 76\% of WIMF solar to offshore wind and reduce solar to 24\% of the WIMF energy share in 2050. 

The WIMF scenario also rely heavily on hydropower, up to 5.5\% of final electricity, more than what is available in the UK. We acknowledge the limitation and keep it as it is.





\subsection{Results}
\label{app_energy_results}


We present the results in terms of energy mix, average price and productivity growth. Productivity growth confidence intervals are the quantiles of the error of forecasting using the value $\zeta=0.126$ on KLEMS database to hindcast productivity growth at the industry level from 1996 to 2021.
We first present the results for the 5 CCC CB6 scenarios, then the 3 WIMF scenarios, a UK adapted WIMF Fast Transition pathway. We then analyse the results in terms of useful energy.

\subsubsection{CCC scenarios}
\label{section_energy_ccc}

Figure \ref{fig_g_cumAhat_CCC} and Table \ref{tab_ccc_cumAhat} shows productivity growth of the energy sector in the CCC scenarios. All scenarios, except one show productivity losses due to the transition. This is due to two main factors: a limited electrification\footnote{Electricity represents 15.6\% of final energy in 2022, In Balanced Net Zero in 2050: 51.7\%, Headwinds: 36.2\%, Tailwinds 58.2\%, Widespread Engagement 60.4\%, Widespread Innovation 58.6\%.} and the use of expensive sources of electricity, such as gas CCGT power plant equiped with CCS.   

Figure \ref{g_scenario_CCC} decomposes productivity gains, plotting the energy mix and the evolution of the average energy prices, and provides three insights: i) We are not exactly comparing the same scenarios, as some scenarios do not reach carbon neutrality, CCS for final energy demand is \textit{not} included because the emissions are decentralised and outside the scope of the energy sector classically defined. Hence, Headwinds remains dependent on fossil fuels. \FL{I don't understand point(i)} ii) Small differences in the energy mix matter. The differences between the  Widespread engagement and Widespread Innovation scenarios is not striking a priori, but it results in 1000bp difference in productivity. iii) electrification is associated with a decrease in the total final demand, as conversion losses from fossil fuel to electricity (about 40\%) are saved. Hence, Headwinds which limits electrification performs better in pure productivity terms. 

\begin{figure}
    \centering
    \caption{Productivity gains in the energy sector according to the CCC scenario}
    \label{fig_g_cumAhat_CCC}
    \includegraphics[width=\textwidth]{figures/figures_energy/g_cumAhat_CCC.jpeg}
     {\\\raggedright\footnotesize\textbf{Note}: Scenarios differ in the evolution of the energy mix, relying on different technologies. Prices are the same in all scenarios. For details of the scenarios refer to figure \ref{g_scenario_CCC}.\\}
\end{figure}

\input{figures/figures_energy/tab_CCC_cumAhat}


\begin{figure}
    \centering
      \caption{Energy mix, average price of energy, cumulative productivity growth of the 5 CCC scenarios for the energy sector}
    \label{g_scenario_CCC}
    \includegraphics[width=0.8\textwidth]{figures/figures_energy/g_scenario_CCC.jpeg}
    {\\\raggedright \footnotesize\textbf{Note}: First column: final energy mix in TWh for the UK, the hatched colours are non-electricity final consumption. CCC 6th Carbon budget scenarios. Second column is the average price of final energy, refer to section \ref{app_energy_prices} for more details. Prices are identical in all scenarios. The dotted line is the average price of energy in 2022, 39.4 USD 2021/MWh. Third column is the cumulative log productivity growth of the sector, $\zeta=0.126$) with intervals of confidence from KLEMS database, representing the error in the approximation of industry-level productivity growth through $\zeta$. Dotted line is no net productivity growth. \\}
\end{figure}


\subsubsection{Variations on CCC scenarios}
\label{section_variation_CCC_wind_solar}
The CCC scenarios incur productivity losses due to their use of expensive electricity generation sources, especially gas and biomass with CCS (4.8\% of total final energy in Balanced Net Zero, 6.1\% in Headwind). We test the impact of these few TWh of electricity generated with expensive technologies: they represent an opportunity cost of up to 2000bp in terms of productivity over 28 years (technically even less, since they are at 0 in 2022, and are developed starting 2026-2027). 

Table \ref{tab_ccc_cumAhat_focus_renewable} shows that if the CCC scenario were to use solar or wind instead of the two mentioned technologies to produce electricity, productivity growth would be significantly impacted. For the Balanced Net Zero scenario, the gain is of 1900bp replacing electricity from gas+CCS, biomass+CCS by Solar PV, from -500 to +1500bp. See Figures \ref{g_scenario_CCC_solar} and \ref{g_scenario_CCC_Wind}. We do not increase the need for storage (P2X or batteries) as the scenarios remain "low renewables".\footnote{ As a sensitivity test, we require that the added quantities of offshore wind ($q_{\text{additional wind}}=q_{\text{elec gas+CCS}}q_{\text{elec biomass+CCS}}$) increase the annual storage by 30\% of the additional \textit{production} $q_{\text{additional storage}}=q_{\text{additional wind}}\times1.5\times0.3$, which is the highest requirement for storage. And that electricity from hydrogen is increased by $q_{\text{additional elec hydrogen}}=q_{\text{additional wind}}\times\frac{1}{12}$. Productivity drops by 60bp compared to the scenario where the transfer of gas+CCS and biomass to CCS to offshore wind without increasing the need for storage.}\footnote{
%
We can cross-check this result using a back of the envelop calculation. In 2050, the costs in 2020 USD are as follows: wind is \$24, biomass+CCS electricity is \$280, and gas+CCS is \$112. In a balanced net-zero scenario, the average price per MWh in 2050 is \$39. In the “wind focus” scenario, it is \$32/MWh. Biomass+CCS electricity and gas+CCS together account for 5\% of the final energy mix in 2050 (1.6\% and 3.3\%, respectively). The difference from \$32 to \$39 can be explained by adding a \$256 premium for biomass and an \$88 premium for gas compared to wind: $32 + 0.01666 \times  (280 - 24) + 0.03333 \times (112 - 24) = 39.2$. A change in price from \$39 to \$32 results in a significant productivity gap: $\hat{p} = \ln(32) - \ln(39) = -0.198$. Productivity change is: $d\hat{A} = (\zeta - 1) \times \hat{p} = (0.126 - 1) \times -0.198 = 0.17 \approx 1730 \, \text{bp} $. The final gap is closer to 1900bp due to the progressive increase in biomass CCS and gas CCS in the energy mix.
%
}


\input{figures/figures_energy/tab_CCC_cumAhat_focus_renewable}

\begin{figure}[H]
    \centering
      \caption{Energy mix, average price of energy, cumulative productivity growth of the "Solar focused Balanced Net Zero scenario"}
    \label{g_scenario_CCC_solar}
    \includegraphics[width=0.9\textwidth]{figures/figures_energy/g_scenario_CCC_solar.jpeg}
    {\\\raggedright \footnotesize\textbf{Note}: See Figure \ref{g_scenario_CCC}. Electricity generated with gas+CCS and biomass+CCS is replaced by solar PV production. We assume no new grid cost nor storage. \\}
\end{figure}

\begin{figure}[H]
    \centering
      \caption{Energy mix, average price of energy, cumulative productivity growth of the ``Wind focused Balanced Net Zero scenario''}
    \label{g_scenario_CCC_Wind}
    \includegraphics[width=0.9\textwidth]{figures/figures_energy/g_scenario_CCC_Wind.jpeg}
    {\\\raggedright \footnotesize\textbf{Note}: See Figure \ref{g_scenario_CCC}. Electricity generated with gas+CCS and biomass+CCS is replaced by offshore wind production. We assume no new grid cost nor storage. \\}
\end{figure}

\subsubsection{Way et al (2022) scenarios}

We implement the global energy sector as forecasted by Way, Ives, Mealy and Farmer  (\citeyear{wayEmpiricallyGroundedTechnology2022}),hereafter the WIMF scenarios.

Table \ref{tab_way_cumAhat} shows in column 4, the cumulative productivity $\mathscr{A}$ over the period 2022-2050. The first column shows the contribution of decrease of the average price of energy to productivity growth, the second column the contribution of the growth in energy demand. Figure \ref{g_scenario_way} shows the energy mix in the three scenario, the average price and cumulative productivity.

First, all three scenario result in a substantial productivity increase. Even the ``No transition scenario'' sees a productivity increase. The average price of energy falls because renewable electricity adds up to the total final energy demand instead of replacing fossil fuel, and because the total energy demand increases to more than 600 EJ/year, compared to the 400 EJ in the Fast transition scenario. 
Second, the fast transition scenario brings the most productivity gains, about 4500bp. A massive electrification (75\% of final energy demand in 2050) limits the increase of production but causes the average price to drop to 26.5 USD 2020/MWh ($\approx$-70\%).
Third, once again, the No Transition scenario does not bring the world to carbon neutrality as it depends massively on fossil fuel in 2050. Extra CCS capacity, most likely DACCS, would have to be deployed. Anyhow, the fast transition stills brings more productivity gains than the No or Slow transition while reducing emissions of the energy sector to net zero. 



\input{figures/figures_energy/tab_cumAhat_way}

\begin{figure}[H]
    \centering
     \caption{Energy mix, average price of energy, cumulative productivity growth of the global energy sector scenarios of \cite{wayEmpiricallyGroundedTechnology2022}}
    \label{g_scenario_way}
    \includegraphics[width=0.9\textwidth]{figures/figures_energy/g_scenario_Way.jpeg}
    {\\\raggedright \footnotesize\textbf{Note}: First column: final energy mix in EJ for the world, the hatched colours are non-electricity final consumption. Second column is the average price of final energy, refer to section \ref{app_energy_prices} for more details. Prices are identical in all scenarios. The dotted line is the average price of energy in 2022, 56.8 USD 2020/MWh. Third column is the cumulative log productivity growth of the sector, $\zeta=0.126$). Dotted line is no net productivity growth.\\}
\end{figure}

\subsubsection{WIMF scenarios applied to the UK}

We describe how we build UK-specific WIMF scenario in section \ref{section_UK_WIMF_scenario}. Productivity gains are +460bp over 2022-2050.\footnote{Without the solar correction, productivity gains would be +1100bp over 2022-2050, as solar is cheaper than wind.}


\begin{figure}[H]
    \centering
     \caption{WIMF scenarios applied to the UK}
    \label{g_scenario_CCC_WIMF}
    \includegraphics[width=0.9\textwidth]{figures/figures_energy/g_scenario_CCC_WIMF.jpeg}
    {\\\raggedright\footnotesize\textbf{Note}: First column: final energy mix in TWh UK, reflecting energy share of the WIMF "Fast Transition" scenario, with a heavier focus on wind than solar. The hatched colours are non-electricity final consumption. Second column is the average price of final energy, refer to section \ref{app_energy_prices} for more details. Prices are identical in all scenarios. The dotted line is the average price of energy in 2022, 39.4 USD 2020/MWh. Third column is the cumulative log productivity growth of the sector, $\zeta=0.126$). Dotted line is no net productivity growth.\\}
\end{figure}


Where does the difference with WIMF global Fast Transition come from ? UK WIMF is 470bp whereas global WIMF is +5300bp. 
The difference lies in three main sources:
\begin{itemize}
    \item 1230bp: Energy production increase (in global) vs decrease in UK (contribution of $\zeta\hat{q}$ in Global WIMF is +485 and -746bp in UK WIMF
    \item 640bp: solar is 2.2 times wind in global WIMF whereas wind is 5.2 times more than solar in CCC balanced net zero (in 2050) (tested with and without the report of 76\% of solar to wind in WIMF scenario). 
    \item 3130: initial price of energy in 2022 in global WIMF is 56.8 USD/MWh whereas it is 39.7USD/MWh in UK WIMF for a similar cost in 2050: 32.5USD/MWh (including storage, network cost, H2). See table \ref{tab_initial_price_energy} for details.
    \begin{itemize}
        \item Global WIMF (0.126-1)*(log(32.5) - log(56.8))= +4880bp
        \item UK WIMF (0.126-1)*(log(32.5) - log(39.7))= +1750bp
    \end{itemize}
\end{itemize}

Summing these three sources, it adds up to 5000bp, thus explaining most of the gap between UK WIMF (470bp) and global WIMF (5300bp). The small difference is due to non linear decrease in price between 2022 and 2050.


\input{figures/figures_energy/tab_initial_price_energy}

\subsubsection{Useful energy - WIMF scenario}

We use the WIMF global scenario for useful energy, instead of final energy. It results in a growing energy demand, between 300 and 400 EJ. Productivity growth is then boosted by the quantity part of the equation, to +7300bp in the Fast Transition scenario. 


\begin{figure}[H]
    \centering
    \caption{Useful energy mix, average price of useful energy, cumulative productivity growth of the global energy sector scenarios of \cite{wayEmpiricallyGroundedTechnology2022}}
    \label{g_scenario_way_useful}
    \includegraphics[width=0.9\textwidth]{figures/figures_energy/g_scenario_way_useful.jpeg}
    {\\\raggedright\footnotesize\textbf{Note}: First column: final energy mix in EJ for the world, the hatched colours are non-electricity final consumption. Second column is the average price of final energy, energy prices are from \cite{wayEmpiricallyGroundedTechnology2022} when available, refer to section \ref{app_energy_prices} for more details. Prices are identical in all scenarios. The dotted line is the average price of energy in 2022,  67.4  USD 2020/MWh. Third column is the cumulative log productivity growth of the sector, $\zeta=0.126$). Dotted line is no net productivity growth.\\}
\end{figure}

\input{figures/figures_energy/tab_cumAhat_way_useful}


\subsubsection{Useful energy - WIMF+CCC scenario}
\label{section_useful_WIMF_UK}
The key sector of the net zero transition is the energy sector, plus the energy savings of the other sectors. One way to represent these two movements is to analyse directly the useful energy sector. We convert final energy into direct energy by using the ratio for each energy source at each year in WIMF scenario and apply them to the final energy WIMF+CCC scenario. Productivity growth is about 1670bp over 2022-2050.

\begin{figure}[H]
    \centering
    \caption{WIMF useful energy scenarios applied to the UK}
    \label{g_scenario_CCC_WIMF_useful}
    \includegraphics[width=0.9\textwidth]{figures/figures_energy/g_scenario_CCC_WIMF_useful.jpeg}
    {\\\raggedright\footnotesize\textbf{Note}: First column: final energy mix in TWh for the world, the hatched colours are non-electricity final consumption. Second column is the average price of final energy, refer to section \ref{app_energy_prices} for more details. Prices are identical in all scenarios. The dotted line is the average price of energy in 2022,  44.1  USD 2020/MWh. Third column is the cumulative log productivity growth of the sector, $\zeta=0.126$). Dotted line is no net productivity growth.\\}
\end{figure}



\subsection{Final or useful energy?}
\label{sect_final_useful_energy}

Should we use useful energy or final energy as the output of the energy sector? We base this section on \cite{wayEmpiricallyGroundedTechnology2022}. Let's create a two-sector world: one energy sector and one manufacturing sector. 

Final energy is ``the energy delivered for use in sectors of the economy'', and useful energy is "the portion of final energy used to provide energy services, such as heat, light, and kinetic energy" \citep{wayEmpiricallyGroundedTechnology2022}. Classic example: an ICE vehicle suffers from conversion losses (converting fuel to kinetic energy) whereas EV do not suffer from the same limitation. An electrification of the pool of cars would then reduce final energy without reducing useful energy.

We consider a two-sector economy: the energy sector and the manufacturing sector. The energy sector transitions from gas to electricity.

\paragraph{Set up}
Let us assume a relationship between a quantity $Q_F$ of gas (fossil fuel) and $Q_G$ of electricity (green) such that they provide the same useful energy: $Q_G=xQ_F$ with $0<x<1$.  $E$ is the energy sector, either green electricity or fossil gas. We consider that there is no change in use, but only a switch from one energy to another so that the useful energy content remains the same. For instance, imagine that the energy sector before the transition produces and sells 10 GWh of gas, and after the transition it sells only 9 GWh of electricity, such that $x=0.9$. The demand of electricity from the manufacturing sector matches the supply of the energy sector (that is, its demand of useful energy is fixed).
  
For the sake of simplicity, let's assume that electricity is slightly more expensive than gas, and that the price of a kWh of electricity is $p_{G}=\frac{1}{x}p_F$ with $p_F$ the price of gas. So the energy bill of the manufacturing sector is constant.




\paragraph{Final energy} 
For the energy sector, production goes does down: 
\begin{equation}
\hat{X}_E=\ln{Q_G}-\ln{Q_F}=\ln{xQ_F}-\ln{Q_F}=\ln{x}
\end{equation}

The price of the product goes up:
\begin{equation}
\hat{p}_E=\ln{p_G}-\ln{p_F}=\ln{\left(\frac{1}{x}p_F\right)}-\ln{p_F}=-\ln{x}
\end{equation}

Hence, productivity growth is
\begin{equation}
\hat{A}_E=\zeta \hat{X}_E+ (\zeta-1)\hat{p}_E=\zeta \ln{x}-(\zeta-1)\ln{x}=\ln{x}. 
\end{equation}

For the manufacturing sector, the quantity of intermediate input goes down while its price goes up. Remember that  $
\hat{A}_i = \hat{X}_i - \hat{\Sigma}_X = \hat{\Sigma}_p - \hat{p}_i$ with the  notation
$
\hat{\Sigma}_X \equiv \sum_j  w_j \hat{X}_{ij}$ and  
$\hat{\Sigma}_p \equiv \sum_j  w_j \hat{p}_{j}$. We recall that we assumed no growth in the production of the manufacturing sector facing a fixed demand.

We then write for the manufacturing sector $C$: 
\begin{equation}
\hat{\Sigma}_X=w_{EC}\hat{X}_{CE}=w_{EC}\ln{x},
\end{equation}
and

\begin{equation}
\hat{\Sigma}_P=w_{EC}\hat{p}_{E}=-w_{EC}\ln{x},
\end{equation}
with $w_{EC}$ the input share of energy in the manufacturing sector: $w_{EC}=\frac{X_{EC}p_E}{X_Cp_C}$. Note that it checks the key assumption $\hat{\Sigma}_X=\frac{1-\zeta}{\zeta}\hat{\Sigma}_P$ when $\zeta=0.5$.

Hence productivity growth is:
\begin{align}
\label{Ahat_C_final_energy}
\hat{A}_C&=\hat{X}_C - \hat{\Sigma}_X = \hat{\Sigma}_p - \hat{p}_C\\
\hat{A}_C&=\hat{X}_C -w_{EC}\ln{x}=-w_{EC}\ln{x}- \hat{p}_C
\end{align}

Now, equation \eqref{Ahat_C_final_energy}gives us $\hat{X}_C=-\hat{p}_C$, and since $\hat{X}_C=0$ by assumption, then $\hat{p}_C=0$ and 
\begin{equation}
\hat{A}_C=-w_{EC}\ln{x}.
\end{equation}
Note that $w_{EC}$ the input share of energy is constant, because the energy expenditures of C is $p_FX_{FC}=\frac{p_G}{x}X_{GC}x=p_GX_{GC}$. The change in energy does not change the energy share of input.

With $\lambda_E$ and $\lambda_C$ the Domar weights of the energy and the manufacturing sectors, aggregate productivity growth is: 
\begin{equation*}
\hat{A}=\lambda_E\ln{x}-\lambda_Cw_{EC}\ln{x}
\end{equation*}

Note that:
\begin{equation*}
    \lambda_Cw_{EC}=\frac{\text{GO}_C}{Y}\frac{X_{EC}p_E}{\text{GO}_C}=\frac{X_{EC}p_E}{Y}=\lambda_E
\end{equation*}

Hence there is no productivity gains
\begin{equation}
\label{eq_Ahat_WE}
\hat{A}=\lambda_E\ln{x}-\lambda_Cw_{EC}\ln{x}=0
\end{equation}

\paragraph{Useful energy} Now let's consider that the energy sector directly sells useful energy. Its production stays the same and its price as well. 
Indeed, the price of a useful kWh of energy before the transition equals the price of useful energy after the transition ($p_G$).
First, the sales of the energy sector are constant, $Q_Gp_G=Q_Fx\frac{p_F}{x}=Q_Fp_F$. Second, the unit price of a kWh of useful energy is


\begin{equation*}
    \frac{\text{Sales}}{\text{Useful energy}}=\frac{Q_F p_F}{xQ_F}=p_F\frac{1}{x}=p_G.
\end{equation*}
Hence, there is no productivity change for the energy sector. 

As for the manufacturing sector, it receives the same quantity of useful energy ($\hat{X}_{EC}=0$) at the same price ($\hat{p}_E=0$). Hence there is no productivity growth either. 

\paragraph{Comparison} Considering the final energy or the useful energy gives the same result: a null aggregate productivity growth. That is excellent news. It is equivalent to consider final energy of useful energy as long as we are considering the full economy. 

\paragraph{Potential for errors}
To decide which metric of energy to use (final or useful), we should take into account the potential for errors for each of the two methods. 

If we use the \textit{useful energy} metric and if by mistake the scenario for the manufacturing sector were to reflect the change in energy supply, then we would have $\hat{A}_C=-w_{EC}\ln{x}$ and then $\hat{A}=-\lambda_Cw_{EC}\ln{x}>0$. We would have overestimated the aggregate productivity growth. This mistake is easily made as we sometimes use aggregate and non transparent cost estimates at the industry-level.

Moreover, note that by using \textit{final energy}  we are missing on the terms $\sum_{i\neq C}\lambda_i w_{Ei}\ln{x}$ for all the industries we are not explicitly considering and we cannot simplify the expression using $\sum_{i\neq C}\lambda_i w_{Ei}=\lambda_E$. The hypothesis here is that the sectors which are not carbon-intensive (e.g. services) are either already using electricity or have a low input share of energy and are thus little affected. In this case as well, we are overestimating the productivity growth. We would not make this error if we were considering useful energy.

\paragraph{Conclusion} It is safer to use \textbf{final energy} rather than useful energy for our computations as the risk is bigger to have energy savings included in the industry-level price estimates than the weight of the sectors not included in the analysis.

